\uuid{3oQ5}
\titre{Somme de variables normales non indépendantes}
\auteur{Maxime Nguyen}
\organisation{AMSCC}
\datecreate{2026-11-27}
\niveau{L2}
\module{Probabilité}
\chapitre{Probabilité continue}
\sousChapitre{Loi normale}
\difficulte{2}
\theme{loi normale, indépendance, somme de variables aléatoires}

\contenu{
\begin{question}
On note $\mathcal{N}(\mu,\sigma)$, pour $\sigma > 0$, la loi normale d'espérance $\mu$ et d'écart-type $\sigma$. Soit $X \sim \mathcal{N}(0, 1)$ et $Y = -X$. Quelles affirmations sont vraies ?
\begin{answers}
\good{$Y$ suit la loi $\mathcal{N}(0,1)$.}
\bad{$X + Y$ suit la loi $\mathcal{N}(0, \sqrt2)$.}
\good{$X + Y$ est la variable constante égale à $0$ ; elle n'a pas de densité.}
\bad{La somme de deux variables de loi normale suit toujours une loi normale.}
\end{answers}
\begin{explanations}
A : la densité de $\mathcal{N}(0,1)$ est paire, donc $-X$ a la même loi que $X$.

B, D sont faux : la règle « les espérances et les variances s'ajoutent, la somme est normale » suppose les variables indépendantes. Ici $X$ et $Y$ ont la même loi que dans le cas indépendant, mais $X + Y = 0$ : une constante n'est pas une variable à densité, donc pas une variable de loi normale.

Contre-sens visé : appliquer la règle de la somme en oubliant
l'hypothèse d'indépendance. Connaître la loi de $X$ et celle de $Y$ ne suffit
pas pour connaître celle de $X + Y$.
\end{explanations}
\end{question}
}
